This element is used for simulating intra-beam scattering (IBS).  The IBS algorithm is based on Bjorken and
Mtingwa~\cite{BM}, and with an extension to include vertical dispersion. It can be used for storage
rings, transport lines, and linacs.

To perform IBS calculations, one or more IBSCATTER elements must be inserted in the beamline. {\tt elegant} calculates
the integrated IBS growth rates between IBSCATTERs (or from beginning of the beamline to the first IBSCATTER), then
scatters particles accordingly at each IBSCATTER element. The beam's parameters are updated for use in downstream
elements.

The IBSCATTER element cannot be can inserted at the beginning of beamline. The number of other elements between
IBSCATTERs or from the beginning of beamline to the first IBSCATTER has to be 2 or more. For storage ring, an IBSCATTER
must be the last element in the beamline.

Because the IBS growth rates are energy dependent, special caution is needed for calculations with accelerating
beam. The user may need to split accelerating cavities into several pieces, so that $\gamma$ has no large changes
between elements.

The user can examine the calculation through an optional SDDS output file.  The file has a multiple page structure. Each
slice at pass $i$ at each IBSCATTER element occupies one page. Each page contains integrated IBS growth rates between
IBSCATTERs (or from the beginning of the beamline to first IBSCATTER) as parameters, and local rates for elements in between
as tabular data.
